\appendix

\chapter{Notation}\label{app:notation}
\begin{center}\footnotesize\setlength{\tabcolsep}{4pt}
\begin{longtable}{@{}p{3.2cm}p{9.8cm}@{}}
\toprule Symbol & Meaning\\\midrule\endhead
$[k]$, $S_k$ & the label set $\{0,\dots,k-1\}$ (or $\{1,\dots,k\}$ where stated) and its permutations\\
$P_k$, $\partial P_k$ & the permutohedron of order $k$ and its boundary complex; faces $=$ ordered set partitions\\
$A\subseteq S_k$ & a set of admissible orderings (histories)\\
$K_A$ & the history complex: faces of $P_k$ all of whose vertices lie in $A$\\
$Y_A$ & the forbidden-history complex, complement of $K_A$ in the sphere of orderings\\
$D_A$, $N[D_A]$ & minimal forbidden prefixes and their nerve (\cref{ch:duality})\\
$G_A$ & legal-swap graph of $A$ with colours (\cref{ch:swaps})\\
$\Pref(w)$ & set of prefix sets of the ordering $w$; $\Pref(F)$ for a face\\
$K(w)$, $K_F$ & thresholds: $\lcm$ of the prefix denominators over $w$, respectively over the vertices of $F$\\
$D(I)$, $\omega_p(I)$ & denominator of a prefix set $I$ and its $p$-adic exponent (\cref{ch:arith})\\
$X_C(m)$ & history complex of the content $C$ at source $m$\\
$T_N(m)$ & the SELECT-II step $m(N+1)^2/N$, $N\mid m$\\
$\beta_j$, $\Hh_j$ & Betti numbers and reduced homology over $\F=\mathbb F_2$ unless stated\\
$f_j$, $\rho_j$ & number of $j$-faces and rank of $\partial_j$ over $\F$\\
$F_t(m_0)$, $A(t)$ & fibre of complete histories to the terminal state $t$ from seed $m_0$; ambiguity $\log_2|F_t|$\\
$\Ent$, $\Mut$ & Shannon entropy and mutual information in bits\\
$\Lambda$, $\first$, $\last$, $L_{c_1,c_2}$ & lift operators of \cref{ch:dimension}\\
$J_X$ & filling-dependence map of \cref{ch:dimension}\\
$B_r(w)$, $I_r$ & ball of radius $r$ in the swap graph and the accessibility profile (\cref{ch:access})\\
$Z_p$, $Z_S$ & subcomplexes of ``adjacency $p$'' in \cref{ch:locality}\\
\bottomrule
\end{longtable}\end{center}

\chapter{Foundational facts}\label{app:lemmas}
These are the classical and elementary facts used without comment in the main text. We state each with a proof or a reference.

\begin{proposition}[faces of the permutohedron]\label{prop:faces}
The faces of $P_k$ are in bijection with ordered set partitions $(B_1|\dots|B_m)$ of $[k]$, the face of the partition having dimension $k-m$ and vertex set the orderings that list $B_1$ first (in any order), then $B_2$, and so on; face $F'$ is a face of $F$ iff the partition of $F'$ refines that of $F$. $P_k$ is simple.
\status{\Classical}{Ziegler, \emph{Lectures on Polytopes}; Bj\"orner--Brenti \cite{Ziegler,BjornerBrenti}}
\end{proposition}

\begin{proposition}[Betti numbers from face counts and ranks]\label{prop:ranknull}
For a finite regular cell complex with boundary matrices $\partial_j$ over a field, $\beta_j=f_j-\rho_j-\rho_{j+1}$, with $\rho_0=0$.
\status{\Analytic}{rank--nullity}
\end{proposition}
\begin{proof}
$\beta_j=\dim\ker\partial_j-\dim\operatorname{im}\partial_{j+1}=(f_j-\rho_j)-\rho_{j+1}$.
\end{proof}

\begin{proposition}[difference formula]\label{prop:diff}
If $X\subseteq X'$ are subcomplexes of one regular cell complex, then for each $j$ one has $\Delta\beta_j=\Delta f_j-\Delta\rho_j-\Delta\rho_{j+1}$, where $\Delta$ denotes $X'$ minus $X$.
\status{\Analytic}{\cref{prop:ranknull} applied twice}
\end{proposition}

\begin{lemma}[gcd and divisibility of thresholds]\label{lem:lcmdiv}
If $F'$ is a face of $F$ then $K_{F'}\mid K_F$, and $F\in X_C(m)$ iff $K_F\mid m$.
\status{\Analytic}{\cref{thm:face-threshold}(a)}
\end{lemma}

\begin{lemma}[cones are contractible; used in the nerve theorem]\label{lem:cone}
A simplicial complex with a vertex adjacent to every vertex of every simplex (a cone) is contractible.
\status{\Classical}{\cite{Hatcher}}
\end{lemma}

\begin{lemma}[Gibbs' inequality]\label{lem:gibbs}
For probability distributions $p,q$ on a finite set, $-\sum_i p_i\log p_i\le-\sum_i p_i\log q_i$ with equality iff $p=q$.
\status{\Classical}{\cite{CoverThomas}}
\end{lemma}

\chapter{Supplementary proofs and discharge status}\label{app:proofs}
The main-text proofs are complete wherever a result is marked analytic. This appendix lists, for every result that is \emph{not} fully proved in the text, exactly what is missing. The same information, organised by theorem, is in the file \texttt{PROOF\_AUDIT.md} of the package.

\section{Results whose proof is incomplete or computer-only}
\begin{center}\footnotesize
\begin{tabular}{@{}p{3.1cm}p{9.9cm}@{}}\toprule Result & What is not proved by hand\\\midrule
\cref{thm:pairwise} (F1), (F2) & The acyclicity of the $30$ single-adjacency complexes and the connectedness of the $435$ pair complexes are certified by exact $\F$-elimination only. Non-emptiness (F2, F3) has a human proof (\cref{lem:Z}). The theorem is for $k=6$ only.\\
\cref{thm:skeleton} & Rests on the classification of $6$-cycles through a vertex of the swap graph, verified by enumeration for $k\le8$; no uniform proof is given.\\
\cref{thm:face-recognition} & Completeness is proved. Soundness for $P_k$-convex sets is computer-assisted for $k\le7$. A proof of soundness for all $k$ (T20d in the source programme) is open.\\
\cref{obs:nondeg} & Non-degeneracy of $554$ off-diagonal pairs is finite-exact; the general statement (LT-6, \cref{prob:nondeg}) is not proved.\\
\cref{thm:JX} & J-1, J-2, J-6 are proved; J-3--J-5 are finite-exact; J-7 (injectivity, \cref{prob:Jinjective}) is open.\\
\cref{obs:anchored} & Injectivity of the history count on the comparable window is a finite-exact statement on two anchors; the general statement (GI-8 of the source programme) is not proved.\\
\cref{prop:gate-nec}, \cref{rem:gate-correction} & The strong form of gate necessity (PS-3 of the source programme) is false (PS-3$'$); only the finite-exact refutation and the exact criterion \cref{thm:exact-survival} are claimed.\\
\bottomrule\end{tabular}\end{center}

\section{Computations recorded as \texttt{NOT\_RERUN}}
The following inherited computations were not re-executed for this monograph; their certificates are named in \cref{app:sources} and their numbers are quoted as inherited: the six-label census of $346{,}859$ source regimes, the seven-label scans, the three-lift search (HL-3), the D21 table of \cref{tab:D21}, the held-out rounds of \cref{obs:atlas-pred}, the atlas rounds of \cref{obs:atlas}, and the values $r=8,9,10$ of the exact history counts for $72^r$.

\chapter{Computational methods}\label{app:methods}
\section{Arithmetic and tolerances}
All computations use exact integer or \texttt{Fraction} arithmetic and $\F$ linear algebra on bit-packed rows. No floating-point value enters a theorem or a table; floating-point appears only in the plotting of the figures, from exact data. The verification engine is the single module \texttt{verify/gigmini.py}, written for this monograph, independent of the engines of the frozen releases.

\section{What was re-derived and what was not}
\Cref{tab:methods} lists the groups of statements by the script that re-derives them, and the result file in which the outcome is recorded.
\begin{table}[ht]\centering\footnotesize
\begin{tabular}{@{}p{3.4cm}p{5.2cm}p{3.7cm}@{}}\toprule Script & Re-derives & Result file\\\midrule
\texttt{run\_\allowbreak checks.py} & histories of $72$, Lemma~F, anchors A and B, the anchor windows, witnesses W1--W6, the radius-two identity, nerve theorem instances, Alexander duality on random sets, the adjacency lemmas F1--F3, the triple counterexample, LT-0, J-2 & \texttt{results.json}\\
\texttt{run\_\allowbreak checks2.py} & the certificate $(z,\varphi)$, the $540$ three-faces, $P_z$ and gates, thresholds, the families $B_0,\dots,B_3$ (PS-5), six-cycles through a root, Petrie versus braid, cube-product property, $72^r$ for $r\le7$, count laws & \texttt{results2.json}\\
\texttt{run\_\allowbreak checks3.py} & the two gate counterexamples (PS-3$'$) & \texttt{results3.json}\\
\texttt{run\_\allowbreak checks4.py} & lift calculus on the anchor: LT-1, LT-2, LT-3, FI-1, FI-3, J-1 & \texttt{results4.json}\\
\texttt{run\_\allowbreak checks5.py} & nerve and core checks on random sets, named systems and all $290$ positives & \texttt{results5.json}\\
\texttt{standalone\_\allowbreak verify.py} (inherited) & the inherited certificates (all $290$ of them): ALL PASS; the tamper test is detected twice & \texttt{standalone\_\allowbreak verify\_\allowbreak manuscript\_\allowbreak run.txt}, \texttt{tamper\_\allowbreak test.txt}\\
\bottomrule\end{tabular}
\caption{Verification scripts and their outputs.}\label{tab:methods}
\end{table}

\section{Figures}
The nine figures are generated by \texttt{figures/make\_figures.py}. Panels whose data are recomputed by the script are those of Figures~\ref{fig:equal-N}, \ref{fig:nerves}, \ref{fig:collapse} (data in \texttt{fig5\_data.json}, \texttt{fig7\_data.json}, \texttt{fig8\_data.json}). Figure~\ref{fig:dimension}(b) and all of \cref{fig:recovery} plot registered inherited values and are not recomputed; panel (c) of \cref{fig:collapse} is hard-coded and was independently reproduced by \texttt{run\_checks5.py}.

\section{Build}
The document is assembled from \texttt{src/} by \texttt{build.sh} and compiled with two \texttt{pdflatex} passes. The quality gates were: zero errors, zero font-substitution warnings, no overfull box above $8$\,pt, no undefined reference or citation. See \texttt{BUILD.md}.

\chapter{Certificate explanations}\label{app:certs}
\section{The anchor certificate}
Let $\lambda_A=(2,3,4,5,8,9)$, with $K_z=288$. A \emph{certificate} for a non-trivial class in $H_2(X;\F)$ is a pair $(z,\varphi)$ with $z$ a $2$-cycle supported in $X$ and $\varphi$ a $2$-cochain with $\varphi\circ\partial_3=0$ on every $3$-face of $P_6$ and $\varphi(z)=1$ (\cref{lem:certificate}). The verification consists of three checks, each independent of the others: $\partial z=0$; $\varphi(\partial F)=0$ for each of the $540$ three-faces $F$ of $P_6$; $\varphi(z)=1$. Because $\varphi$ is defined on $P_6$ and not on $X$, it vanishes on the boundaries of all $3$-faces of any subcomplex: $[z]\ne0$ in every $X$ containing $z$, which is how the survival results are used.

\subsection*{The cycle $z$ ($96$ two-faces)}
Each face is written as an ordered set partition of $\{0,\dots,5\}$ into four blocks (blocks separated by \texttt{|}); block sizes $(3,1,1,1)$ in some order give a hexagon and $(2,2,1,1)$ a square. There are $64$ squares and $32$ hexagons. The list is sorted lexicographically.
\begin{center}\footnotesize\setlength{\tabcolsep}{6pt}
\begin{tabular}{llll}
\texttt{012|3|4|5} & \texttt{014|2|3|5} & \texttt{01|24|3|5} & \texttt{01|2|34|5} \\
\texttt{025|1|3|4} & \texttt{02|15|3|4} & \texttt{02|1|35|4} & \texttt{02|1|3|45} \\
\texttt{045|1|2|3} & \texttt{04|15|2|3} & \texttt{04|1|25|3} & \texttt{04|1|2|35} \\
\texttt{05|12|3|4} & \texttt{05|14|2|3} & \texttt{05|1|24|3} & \texttt{05|1|2|34} \\
\texttt{0|125|3|4} & \texttt{0|12|35|4} & \texttt{0|12|3|45} & \texttt{0|145|2|3} \\
\texttt{0|14|25|3} & \texttt{0|14|2|35} & \texttt{0|15|24|3} & \texttt{0|15|2|34} \\
\texttt{0|1|245|3} & \texttt{0|1|24|35} & \texttt{0|1|25|34} & \texttt{0|1|2|345} \\
\texttt{124|0|3|5} & \texttt{12|04|3|5} & \texttt{12|0|34|5} & \texttt{14|02|3|5} \\
\texttt{1|024|3|5} & \texttt{1|02|34|5} & \texttt{245|0|1|3} & \texttt{24|15|0|3} \\
\texttt{24|1|05|3} & \texttt{24|1|0|35} & \texttt{24|5|01|3} & \texttt{25|04|1|3} \\
\texttt{25|0|14|3} & \texttt{25|0|1|34} & \texttt{2|013|4|5} & \texttt{2|035|1|4} \\
\texttt{2|03|15|4} & \texttt{2|03|1|45} & \texttt{2|05|13|4} & \texttt{2|0|135|4} \\
\texttt{2|0|13|45} & \texttt{2|134|0|5} & \texttt{2|13|04|5} & \texttt{2|14|03|5} \\
\texttt{2|1|034|5} & \texttt{2|345|0|1} & \texttt{2|34|01|5} & \texttt{2|34|05|1} \\
\texttt{2|34|0|15} & \texttt{2|35|04|1} & \texttt{2|35|0|14} & \texttt{2|3|014|5} \\
\texttt{2|3|045|1} & \texttt{2|3|04|15} & \texttt{2|3|05|14} & \texttt{2|3|0|145} \\
\texttt{2|45|03|1} & \texttt{2|45|0|13} & \texttt{2|4|135|0} & \texttt{2|4|13|05} \\
\texttt{2|4|15|03} & \texttt{2|4|1|035} & \texttt{2|4|35|01} & \texttt{2|4|3|015} \\
\texttt{2|4|5|013} & \texttt{2|5|034|1} & \texttt{2|5|03|14} & \texttt{2|5|04|13} \\
\texttt{2|5|0|134} & \texttt{45|02|1|3} & \texttt{45|0|12|3} & \texttt{4|015|2|3} \\
\texttt{4|01|25|3} & \texttt{4|01|2|35} & \texttt{4|125|0|3} & \texttt{4|12|05|3} \\
\texttt{4|12|0|35} & \texttt{4|15|02|3} & \texttt{4|1|025|3} & \texttt{4|1|02|35} \\
\texttt{4|25|01|3} & \texttt{4|5|012|3} & \texttt{5|024|1|3} & \texttt{5|02|14|3} \\
\texttt{5|02|1|34} & \texttt{5|04|12|3} & \texttt{5|0|124|3} & \texttt{5|0|12|34} \\
\end{tabular}\end{center}

\subsection*{The cochain $\varphi$ (six hexagons)}
\texttt{245|0|1|3}, \texttt{245|0|3|1}, \texttt{245|1|0|3}, \texttt{245|1|3|0}, \texttt{245|3|0|1}, \texttt{245|3|1|0}.

\subsection*{The set EXTRA ($20$ prefix sets) and the minimal families}
Prefix sets are written as strings of the labels (ranks) they contain. EXTRA $=\{3, 03, 13, 013, 024, 34, 034, 134, 0134, 15, 125, 35, 035, 135, 0135, 1235, 345, 0345, 1345, 01345\}$.
\begin{itemize}
\item $B_0$ (1 sets): 024.
\item $B_1$ (3 sets): 15, 125, 1235.
\item $B_2$ (15 sets): 3, 03, 13, 013, 34, 034, 134, 0134, 35, 035, 0135, 345, 0345, 1345, 01345.
\item $B_3$ (17 sets): 3, 03, 13, 013, 34, 134, 0134, 15, 35, 035, 135, 0135, 1235, 345, 0345, 1345, 01345.
\end{itemize}


\section{Reading the certificates}
The Boolean statement in \cref{thm:exact-survival} says that $z$ bounds in $X_\lambda(m)$ with $K_z\mid m$ iff the family $S=\mathfrak F_\lambda(m)\cap\mathrm{EXTRA}$ contains one of $B_0,\dots,B_3$. Each $B_i$ was checked to bound and to be minimal under removal of one set (\texttt{results2.json}, \texttt{ch6\_PS5\_minimal\_families}); the criterion was tested against exact filling on $400$ random subfamilies of the $20$ sets and agreed in all $400$. The count $589{,}839$ of the $2^{20}$ families which bound is inherited and was not recounted.

\section{Inherited certificate files}
The files in \texttt{verify/inherited/} are verbatim copies of the frozen certificates used by \texttt{standalone\_verify.py}. Their sha256 hashes begin: \texttt{standalone\_verify.py} 0e683b4f, \texttt{positives\_cycle\_cocycle.json} 88681b94, \texttt{pattern\_H2.json} 58e21df3, \texttt{H2\_anchor\_A.json} a44490e4; full hashes are in \texttt{MANIFEST.json}.

\chapter{Research-source index}\label{app:sources}
\section{Release folders and manifest hashes}
The mathematics of this monograph was integrated from the following frozen releases. The hash is the first sixteen hexadecimal digits of the sha256 of the release's \texttt{MANIFEST.json} as present in the build workspace.
\begin{center}\footnotesize\setlength{\tabcolsep}{5pt}
\begin{tabular}{@{}llll@{}}\toprule Release & Hash & Release & Hash\\\midrule
GIG-I v001 & 28385614d8999629 & GIG-INFO v001 & cb03261b27c3ffc9\\
GIG-I v002 & 0212b57ee3306c6f & GIG-INFO v002 & 7709cc7a02d31fa4\\
GIG-I v003 & fe59b68f8dcbb0b8 & GIG-INFO v003 & 6e79f1013cca47bb\\
GIG-I v004 & 26e14d978aafe3ff & GIG-INFO v004 & abf47ee29f59bebe\\
GIG-I v005 & aea157c8fd126a00 & GIG-INFO v005 & c01a7298716615d0\\
GIG-I v005 S2 & 9558b1071737ff35 & GIG-INFO v006 & 9e4c14ede3e65fe6\\
GIG-I v005 S3 & e38b5e73d4869902 & & \\
GIG-I v006 & ac17c9efddd9f692 & & \\
GIG-I v007 & 2d1d3aa0094c4a8c & & \\
\bottomrule\end{tabular}\end{center}
The frozen nine-page paper \emph{Forbidden Histories and the Topology of Admissible Generation} (version 1) feeds \cref{ch:duality,ch:minimal,ch:locality}.

\section{Erratum on parent hashes}
The fifth GIG-INFO release recorded the same parent hash for GIG-I v004 and v005, because a global substitution in a path helper rewrote only the first occurrence on one line. The true GIG-I v004 manifest hash is the one listed above. The sixth release re-tested a one-token repair of the fifth: all mathematical tests passed and the affected output regenerated identically; no result of the fifth release depends on a module of GIG-I v004 alone. The parent manifests recorded in the sixth release are those present in the build workspace on 2026-11-06; if an earlier archive of a parent has a different hash, the two are different builds, and the point should be settled by comparison with the archives on the user's device. In the build of this monograph $11$ of $14$ recorded parent manifests were matched to the workspace copies; the other $3$ did not match, or could not be matched, and remain unresolved.

\section{Repository links}
No pinned commit of the project repository was available when this monograph was prepared. Consequently the bibliography and this index do not print repository links, and none have been invented. Once a commit is fixed, a link to each release folder at that commit can be added to this table.

\chapter{Theorem crosswalk (summary)}\label{app:crosswalk}
The complete crosswalk, mapping every theorem and number of the monograph to its source release, status and evidence, is the file \texttt{THEOREM\_CROSSWALK.md}. The principal results are summarised here.
\begin{center}\footnotesize\setlength{\tabcolsep}{4pt}
\begin{longtable}{@{}p{3.6cm}p{4.2cm}p{4.8cm}@{}}
\toprule Result & Content & Status\\\midrule\endhead
\cref{thm:closure}, \cref{thm:block} & closure and ordered-block factorisation & analytic\\
\cref{thm:face-threshold} & face thresholds, filtration & analytic\\
\cref{thm:swap-reconstruction} & complex from coloured swap graph & analytic\\
\cref{thm:skeleton} & $2$-skeleton from uncoloured graph, $k\le8$ & analytic modulo finite enumeration\\
\cref{thm:hole-survival}, \cref{thm:anchor} & hole survival, anchor & analytic / finite-exact\\
\cref{thm:exact-survival} & exact survival criterion on $P_6$ & finite-exact\\
\cref{thm:AT1} & survival law for a fixed cycle & analytic\\
\cref{thm:lift-cycle}, \cref{thm:cochain-lift} & lift calculus & analytic / finite-exact\\
\cref{thm:JX} & filling-dependence map & analytic (J-1, J-2, J-6); finite-exact; open (J-7)\\
\cref{thm:obs-ladder}, \cref{thm:profile}, \cref{thm:radius2} & observation ladder, accessibility, radius two & analytic\\
\cref{thm:AD}, \cref{thm:nerve} & duality, dead-set nerve & analytic given classical theorems\\
\cref{thm:barmak}, \cref{cor:core-suffices} & cores & classical, corollary analytic\\
\cref{thm:pairwise} & pairwise locality, $k=6$ & computer-assisted\\
\bottomrule\end{longtable}\end{center}
